\documentclass[12pt]{article}
\usepackage{ficha}

\begin{document}

\ficha{Números Racionais}[Frações e Decimais com Sinal]

\instrucoes{Resolve sem calculadora. Apresenta as frações na forma irredutível.}

\secao{Representação e Comparação}{Responde a cada questão.}
\begin{exercicios}
  \item Indica o simétrico de $-\dfrac{3}{4}$ e de $2{,}5$.
  \item Calcula $\left|-\dfrac{5}{6}\right|$ e $|-0{,}4|$.
  \item Escreve $-\dfrac{3}{5}$ na forma de dízima.
  \item Escreve $-1{,}25$ na forma de fração irredutível.
  \item Compara: $\;-\dfrac{1}{2} \;\square\; -\dfrac{1}{3}$
  \item Compara: $\;-0{,}7 \;\square\; -\dfrac{3}{4}$
  \item Compara: $\;-2{,}1 \;\square\; -2{,}01$
  \item Ordena por ordem crescente: $-1{,}5,\ \dfrac{1}{4},\ -\dfrac{2}{3},\ 0,\ -0{,}2$.
  \item Representa numa reta numérica: $-\dfrac{3}{2},\ -0{,}5,\ \dfrac{3}{4},\ 1{,}25$.
  \item Indica um número racional entre $-\dfrac{1}{2}$ e $-\dfrac{1}{4}$.
\end{exercicios}

\secao{Adição e Subtração}{Calcula.}
\begin{exercicios}
  \item $-\dfrac{1}{4} + \dfrac{3}{4}$
  \item $-\dfrac{2}{3} - \dfrac{1}{3}$
  \item $\dfrac{1}{2} - \dfrac{5}{6}$
  \item $-\dfrac{3}{5} + \left(-\dfrac{1}{2}\right)$
  \item $-\dfrac{3}{4} - \left(-\dfrac{5}{8}\right)$
  \item $-1{,}5 + 0{,}75$
  \item $-2{,}4 - 1{,}35$
  \item $0{,}5 - \dfrac{3}{4}$
  \item $-1 + \dfrac{2}{5} - 0{,}3$
  \item $\dfrac{1}{3} - \left(\dfrac{1}{2} - \dfrac{5}{6}\right)$
\end{exercicios}

\secao{Multiplicação e Divisão}{Calcula.}
\begin{exercicios}
  \item $-\dfrac{2}{3} \times \dfrac{9}{4}$
  \item $\left(-\dfrac{5}{6}\right) \times \left(-\dfrac{3}{10}\right)$
  \item $-4 \times \dfrac{3}{8}$
  \item $-\dfrac{4}{5} : \dfrac{2}{5}$
  \item $\dfrac{3}{7} : \left(-\dfrac{9}{14}\right)$
  \item $\left(-\dfrac{1}{2}\right) : (-4)$
  \item $-0{,}6 \times 1{,}5$
  \item $-2{,}4 : (-0{,}8)$
  \item $\left(-\dfrac{2}{3}\right)^2$
  \item $\left(-\dfrac{1}{2}\right)^3$
\end{exercicios}

\secao{Expressões Numéricas}{Calcula o valor de cada expressão.}
\begin{exercicios}
  \item $-\dfrac{1}{2} + \dfrac{1}{3} \times \dfrac{3}{4}$
  \item $\left(-\dfrac{1}{2} + \dfrac{1}{3}\right) \times 6$
  \item $2 - \dfrac{3}{4} : \left(-\dfrac{1}{2}\right)$
  \item $\left(\dfrac{2}{5} - 1\right)^2$
  \item $-\dfrac{3}{2} \times \left(\dfrac{1}{3} - \dfrac{5}{6}\right)$
  \item $\dfrac{-\frac{1}{2} + \frac{3}{4}}{\frac{1}{8}}$
  \item $-0{,}5 \times \left(2 - \dfrac{1}{5}\right) + 1$
  \item $\left(-\dfrac{2}{3}\right)^2 - \dfrac{1}{9} : \left(-\dfrac{1}{3}\right)$
\end{exercicios}

\secao{Problemas}{Resolve.}
\begin{exercicios}
  \item Às 6h a temperatura era $-2{,}5$\,°C. Até ao meio-dia subiu $7{,}8$\,°C. Qual era a temperatura ao meio-dia?
  \item Um mergulhador está a $-12{,}6$\,m. Sobe $\dfrac{1}{3}$ dessa profundidade. A que profundidade fica?
  \item Uma ação na bolsa valia 24 euros. Na segunda-feira variou $-\dfrac{1}{8}$ do valor e na terça-feira $+1{,}5$ euros. Quanto vale agora?
  \item O saldo de uma conta era $-35{,}40$ euros. Depositaram-se 60 euros e pagou-se uma conta de 18,75 euros. Qual é o saldo final?
  \item A média de cinco temperaturas foi $-1{,}2$\,°C. Qual é a soma das temperaturas?
\end{exercicios}

\secao{Desafios}{Resolve os problemas.}
\begin{exercicios}
  \item Determina $x$ tal que $x + \dfrac{3}{4} = -\dfrac{1}{2}$.
  \item Determina $x$ tal que $-\dfrac{2}{3}\,x = \dfrac{4}{9}$.
  \item O inverso de $-\dfrac{5}{3}$ é? E o inverso de $-0{,}25$?
  \item Indica três números racionais entre $-0{,}1$ e $0$.
  \item Calcula $\left(1 - \dfrac{1}{2}\right)\left(1 - \dfrac{1}{3}\right)\left(1 - \dfrac{1}{4}\right)\cdots\left(1 - \dfrac{1}{10}\right)$.
  \item Se $a = -\dfrac{1}{2}$ e $b = \dfrac{2}{3}$, calcula $a^2 - b$, $\;a \times b$ e $\;\dfrac{a}{b}$.
  \item Explica por que razão o produto de um número negativo pelo seu inverso é sempre $1$.
\end{exercicios}

\end{document}
